
\documentclass[11pt,a4paper]{article}

\usepackage[ngerman]{babel}
\usepackage[T1]{fontenc}
\usepackage[utf8]{inputenc}
\usepackage{geometry}
\usepackage{enumitem}
\usepackage{longtable}
\usepackage{array}
\usepackage{hyperref}
\usepackage{setspace}
\usepackage{titlesec}

\geometry{left=25mm,right=25mm,top=25mm,bottom=25mm}
\setstretch{1.15}

\setlist[itemize]{noitemsep, topsep=3pt}
\setlist[enumerate]{noitemsep, topsep=3pt}

\titleformat{\section}{\large\bfseries}{\thesection.}{0.5em}{}
\titleformat{\subsection}{\normalsize\bfseries}{\thesubsection.}{0.5em}{}
\titleformat{\subsubsection}{\normalsize\bfseries}{\thesubsubsection.}{0.5em}{}

\title{
    \textbf{Steuerliche Dokumentation}\\
    \large DAC6-Negativprüfung, Fremdvergleichsdokumentation und vGA-Prüfung\\
    \vspace{0.5em}
    \normalsize Grenzüberschreitende Dienstleistungsbeziehung Deutschland -- Bulgarien
}

\author{}
\date{\today}

\begin{document}

\maketitle

\section*{Dokumentationszweck}

Dieses Dokument dient der internen steuerlichen Dokumentation einer grenzüberschreitenden Dienstleistungsbeziehung zwischen einer deutschen GmbH und einem in Bulgarien steuerlich ansässigen selbständigen Leistungserbringer. Es umfasst:

\begin{enumerate}
    \item eine DAC6-Negativprüfung
    \item eine Fremdvergleichsdokumentation,
    \item eine Prüfung möglicher verdeckter Gewinnausschüttungen (vGA).
\end{enumerate}

Die Dokumentation dient der nachvollziehbaren Einordnung des Sachverhalts für interne Zwecke, für die laufende steuerliche Dokumentation sowie gegebenenfalls für die Vorlage gegenüber dem steuerlichen Berater oder im Rahmen einer steuerlichen Außenprüfung.

Die Beurteilung erfolgt auf Grundlage des nachfolgend dargestellten Sachverhalts. Bei wesentlichen Änderungen des Sachverhalts, insbesondere hinsichtlich Funktionsumfang, Vergütungsstruktur, Leistungsinhalt, gesellschaftsrechtlicher Stellung oder tatsächlicher Durchführung, ist eine erneute steuerliche Würdigung vorzunehmen.

\section{Beteiligte Parteien}

\subsection{Gesellschaft}

\textbf{TEGRITY Safety Engineering GmbH}\\
Schloßlände 26\\
85049 Ingolstadt\\
Deutschland

\subsection{Gesellschafter und selbständiger Leistungserbringer}

\textbf{Stefan Koch}\\
Ulitsa Georgi Dimitrov 19A\\
4176 Matenitsa\\
Bulgarien

\vspace{0.5em}

Stefan Koch ist Gesellschafter der TEGRITY Safety Engineering GmbH. Er ist nach dem hier zugrunde gelegten Sachverhalt nicht Geschäftsführer der GmbH und übt keine organschaftlichen oder geschäftsführenden Funktionen für die GmbH aus. Die Leistungserbringung erfolgt als selbständiger Berater auf Grundlage eines Beratervertrags beziehungsweise einzelner Workpackages.

\section{Sachverhalt}

Die TEGRITY Safety Engineering GmbH bezieht fachlich abgegrenzte Beratungs- und Projektleistungen von Stefan Koch als in Bulgarien ansässigem selbständigem Leistungserbringer.

Die Leistungen werden auf Grundlage eines Beratervertrags und/oder einzelner projektbezogener Workpackages erbracht. Die Workpackages betreffen Dienstleistungen, die der operativen Leistungserbringung der GmbH gegenüber deren Endkunden dienen. Die GmbH bleibt Vertragspartnerin des jeweiligen Endkunden und trägt weiterhin die unternehmerische Verantwortung gegenüber dem Endkunden.

Der Leistungserbringer:

\begin{itemize}
    \item ist Gesellschafter der GmbH,
    \item ist nicht Geschäftsführer der GmbH,
    \item übernimmt keine organschaftlichen Aufgaben,
    \item erbringt fachlich abgegrenzte Beratungs- oder Projektleistungen,
    \item rechnet projektbezogen, workpackagebezogen oder nach einer vergleichbaren sachlichen Vergütungsgrundlage ab.
\end{itemize}

Die Leistungen:

\begin{itemize}
    \item werden tatsächlich erbracht,
    \item sind fachlich und inhaltlich abgrenzbar,
    \item werden anhand von Workpackages, Leistungsbeschreibungen, Arbeitsergebnissen, Projektdokumentation, Rechnungen oder sonstigen geeigneten Unterlagen dokumentiert,
    \item dienen der operativen Leistungserbringung der GmbH gegenüber Endkunden,
    \item werden nicht allein aufgrund der Gesellschafterstellung vergütet.
\end{itemize}

Es erfolgt keine Übertragung, Nutzung oder Lizenzierung immaterieller Wirtschaftsgüter. Es werden keine Geschäftsführungsfunktionen, keine organschaftlichen Aufgaben und keine gesellschaftsrechtlichen Sonderrechte vergütet.

\section{DAC6-Negativprüfung}

\subsection{Prüfungsgegenstand}

Zu prüfen ist, ob die dargestellte Struktur eine meldepflichtige grenzüberschreitende Steuergestaltung im Sinne der Richtlinie (EU) 2018/822 (DAC6) darstellt.

Eine DAC6-Meldepflicht setzt nicht allein einen grenzüberschreitenden Bezug voraus. Erforderlich ist zusätzlich, dass mindestens eines der gesetzlich definierten Kennzeichen (Hallmarks) erfüllt ist. Bei bestimmten Kennzeichen ist außerdem erforderlich, dass der sogenannte Main-Benefit-Test erfüllt ist, also der Hauptvorteil oder einer der Hauptvorteile der Gestaltung in der Erlangung eines steuerlichen Vorteils liegt.

\subsection{Grenzüberschreitender Sachverhalt}

Ein grenzüberschreitender Sachverhalt liegt vor, da:

\begin{itemize}
    \item der Leistungsempfänger, die TEGRITY Safety Engineering GmbH, in Deutschland ansässig ist,
    \item der Leistungserbringer, Stefan Koch, in Bulgarien ansässig ist.
\end{itemize}

Damit ist eine Prüfung nach DAC6 grundsätzlich angezeigt. Der grenzüberschreitende Bezug allein begründet jedoch keine Meldepflicht.

\subsection{Prüfung der Hallmarks gemäß Anhang IV DAC6}

\subsubsection*{Kategorie A -- Allgemeine Kennzeichen}

Die Voraussetzungen der Kategorie A sind nicht erfüllt.

Insbesondere:

\begin{itemize}
    \item Es bestehen keine Vertraulichkeitsvereinbarungen hinsichtlich steuerlicher Vorteile.
    \item Die Vergütung ist nicht erfolgsabhängig im Hinblick auf einen steuerlichen Vorteil.
    \item Es liegt kein standardisiertes oder vermarktetes Steuermodell vor.
    \item Die Leistungsbeziehung beruht auf einer konkreten operativen Dienstleistung und nicht auf einem serienmäßig angebotenen steuerlichen Strukturierungsmodell.
\end{itemize}

\subsubsection*{Kategorie B -- Spezifische Kennzeichen}

Die Voraussetzungen der Kategorie B sind nicht erfüllt.

Insbesondere:

\begin{itemize}
    \item Es erfolgt keine Umwandlung von Einkünften in eine steuerlich günstiger behandelte Einkunftsart.
    \item Es wird kein bestehendes Arbeitsverhältnis ersetzt.
    \item Es liegt keine Struktur vor, die darauf gerichtet ist, steuerlich höher belastetes Einkommen in niedriger besteuertes Einkommen umzuwandeln.
    \item Es werden keine Tätigkeiten vergütet, die typischerweise einem Geschäftsführer oder Organ der GmbH zuzuordnen wären.
    \item Die Leistungserbringung erfolgt auf Grundlage einer selbständigen, projektbezogenen Dienstleistungsbeziehung.
\end{itemize}

Die Zahlungen an den Leistungserbringer erfolgen für tatsächlich erbrachte Beratungs- oder Projektleistungen. Sie erfolgen nicht als Ersatz für Arbeitslohn, Geschäftsführervergütung oder gesellschaftsrechtliche Gewinnverwendung.

\subsubsection*{Kategorie C -- Grenzüberschreitende Zahlungen}

Die Voraussetzungen der Kategorie C sind nicht erfüllt.

Insbesondere:

\begin{itemize}
    \item Bulgarien ist kein Staat mit keiner oder nahezu keiner Besteuerung.
    \item Bulgarien ist kein nicht kooperatives Steuerhoheitsgebiet im Sinne einschlägiger Listen.
    \item Es liegt kein Präferenz- oder Sondersteuerregime ohne effektive Besteuerung vor.
    \item Es bestehen keine doppelten Abzugs- oder Befreiungseffekte.
    \item Es handelt sich nicht um Lizenzzahlungen, Zinszahlungen oder vergleichbare passive Einkünfte, sondern um Vergütung für operative Dienstleistungen.
\end{itemize}

Die Zahlung beruht auf einem tatsächlichen Leistungsaustausch zwischen der GmbH und dem selbständigen Leistungserbringer. Der Umstand, dass der Leistungserbringer in Bulgarien ansässig ist, begründet für sich genommen kein Kennzeichen im Sinne der Kategorie C.

\subsubsection*{Kategorie D -- Umgehung von Transparenzpflichten}

Die Voraussetzungen der Kategorie D sind nicht erfüllt.

Insbesondere:

\begin{itemize}
    \item Es erfolgt keine Verschleierung wirtschaftlich Berechtigter.
    \item Die Gesellschafterstellung von Stefan Koch ist offen.
    \item Die Leistungsbeziehung ist vertraglich und buchhalterisch dokumentiert.
    \item Es werden keine Melde- oder Transparenzpflichten umgangen.
\end{itemize}

\subsubsection*{Kategorie E -- Verrechnungspreise}

Die Voraussetzungen der Kategorie E sind nicht erfüllt.

Insbesondere:

\begin{itemize}
    \item Es werden keine immateriellen Wirtschaftsgüter übertragen.
    \item Es liegen keine schwer bewertbaren immateriellen Wirtschaftsgüter vor.
    \item Es erfolgt keine konzerninterne Restrukturierung.
    \item Es werden keine Funktionen, Risiken oder Wirtschaftsgüter in einer Weise verlagert, die einen der relevanten Verrechnungspreis-Hallmarks auslöst.
    \item Gegenstand der Leistungsbeziehung sind operative Dienstleistungen.
\end{itemize}

\subsection{Main-Benefit-Test}

Soweit einzelne Hallmarks einen Main-Benefit-Test voraussetzen, ist dieser nicht erfüllt.

Der Hauptzweck der Struktur liegt in der tatsächlichen Erbringung operativer Dienstleistungen für Kundenprojekte der GmbH. Die Leistungen sind wirtschaftlich erforderlich, tatsächlich ausgeführt und für die Leistungserbringung der GmbH gegenüber ihren Endkunden nutzbar.

Ein steuerlicher Vorteil stellt weder den Hauptzweck noch einen der Hauptzwecke der Gestaltung dar. Steuerliche Effekte ergeben sich lediglich als Folge der Ansässigkeit der Beteiligten in unterschiedlichen Staaten und nicht als Ziel der Struktur.

Für die Einordnung spricht insbesondere:

\begin{itemize}
    \item Die GmbH benötigt die Leistungen zur Erfüllung eigener Kundenprojekte.
    \item Der Leistungserbringer verfügt über die für die Leistungserbringung erforderliche fachliche Qualifikation.
    \item Die Leistungen werden tatsächlich erbracht und dokumentiert.
    \item Die Vergütung soll dem Fremdvergleich entsprechen.
    \item Die Struktur ist nicht auf die Erlangung eines steuerlichen Vorteils ausgerichtet.
\end{itemize}

\subsection{Abgrenzung zu steuerlichem Missbrauch}

Es bestehen keine Anhaltspunkte für eine missbräuchliche Gestaltung. Die gewählte Struktur ist wirtschaftlich begründet, da tatsächlich benötigte fachliche Leistungen erbracht werden. Es liegt keine künstliche oder unangemessene Gestaltung vor, die auf die Erlangung eines unangemessenen steuerlichen Vorteils gerichtet ist.

\subsection{DAC6-Ergebnis}

Zusammenfassend ergibt sich:

\begin{itemize}
    \item Grenzüberschreitender Sachverhalt: ja.
    \item Hallmarks gemäß Anhang IV DAC6: nein.
    \item Main-Benefit-Test: nicht erfüllt.
\end{itemize}

Die vorliegende Struktur stellt daher keine meldepflichtige grenzüberschreitende Steuergestaltung im Sinne der DAC6 dar. Eine DAC6-Meldung ist auf Grundlage des dargestellten Sachverhalts nicht erforderlich.

\section{Fremdvergleichsdokumentation}

\subsection{Zweck der Fremdvergleichsdokumentation}

Zweck dieser Dokumentation ist die Darlegung, dass die Vergütung des in Bulgarien ansässigen selbständigen Leistungserbringers dem entspricht, was unabhängige Dritte unter vergleichbaren Umständen vereinbart hätten.

Da Stefan Koch Gesellschafter der TEGRITY Safety Engineering GmbH ist, ist eine nachvollziehbare Dokumentation der Fremdüblichkeit sinnvoll. Maßgeblich ist, ob Leistung, Preis, Vertragsgestaltung und tatsächliche Durchführung einem Drittvergleich standhalten.

\subsection{Leistungsbeschreibung}

Die Leistungen werden auf Grundlage eines Beratervertrags und/oder einzelner Workpackages erbracht. Die Workpackages sollten jeweils mindestens folgende Informationen enthalten:

\begin{itemize}
    \item Projektbezeichnung,
    \item Leistungsgegenstand,
    \item Zeitraum der Leistungserbringung,
    \item konkrete Deliverables oder Arbeitsergebnisse,
    \item Abnahmekriterien oder Nachweisform,
    \item Vergütungsgrundlage.
\end{itemize}

Die Leistungen sind fachlich abgrenzbar und dienen unmittelbar oder mittelbar der Erfüllung von Kundenprojekten der GmbH.

\subsection{Funktions- und Risikoanalyse}

\begin{longtable}{|p{0.32\textwidth}|p{0.58\textwidth}|}
\hline
\textbf{Kriterium} & \textbf{Einordnung} \\
\hline
Funktion des Leistungserbringers & Selbständige Erbringung fachlich abgegrenzter Beratungs- oder Projektleistungen. Keine Geschäftsführung, keine organschaftliche Funktion. \\
\hline
Funktion der GmbH & Vertragspartnerin des Endkunden, Projektverantwortung, Kundenbeziehung, Abrechnung gegenüber Endkunden, wirtschaftliches Unternehmerrisiko. \\
\hline
Risiko des Leistungserbringers & Risiko ordnungsgemäßer Leistungserbringung im Rahmen des Beratervertrags; keine Übernahme des Endkundenrisikos der GmbH, soweit nicht gesondert vereinbart. \\
\hline
Risiko der GmbH & Kundenrisiko, Projektmarge, Vertragserfüllung gegenüber Endkunden, wirtschaftliche Gesamtverantwortung. \\
\hline
Wirtschaftsgüter & Keine Übertragung oder Lizenzierung immaterieller Wirtschaftsgüter. Nutzung eigener Fachkenntnisse des Leistungserbringers zur Dienstleistungserbringung. \\
\hline
\end{longtable}

\subsection{Vergütungsmethode}

Für operative Beratungs- und Projektleistungen ist grundsätzlich eine kosten-, zeit- oder werkbezogene Vergütung fremdüblich. In Betracht kommen insbesondere:

\begin{itemize}
    \item Tagessatz,
    \item Stundensatz,
    \item Festpreis je Workpackage,
    \item Meilensteinvergütung.
\end{itemize}

Die gewählte Vergütungsmethode ist fremdüblich, wenn sie bei vergleichbaren externen Dienstleistern ebenfalls angewendet werden könnte und die GmbH bei objektiver Betrachtung bereit wäre, denselben Preis auch an einen unabhängigen Dritten zu zahlen.

\subsection{Fremdvergleichskriterien}

Für die Beurteilung der Fremdüblichkeit sind insbesondere folgende Kriterien relevant:

\begin{itemize}
    \item Qualifikation und Erfahrung des Leistungserbringers,
    \item Komplexität der Leistung,
    \item Marktüblichkeit vergleichbarer Beratungsleistungen,
    \item Projektverantwortung und fachlicher Beitrag,
    \item Laufzeit und Planbarkeit der Beauftragung,
    \item Dokumentation der erbrachten Leistungen,
    \item Verhältnis zwischen Einkaufspreis der Leistung und Endkundenvergütung der GmbH,
    \item wirtschaftliche Marge der GmbH im Endkundenprojekt.
\end{itemize}

\subsection{Mindestanforderungen an Nachweise}

Zur Absicherung der Fremdüblichkeit sollten folgende Unterlagen vorgehalten werden:

\begin{itemize}
    \item Beratervertrag,
    \item Workpackage-Beschreibungen,
    \item Rechnungen,
    \item Leistungsnachweise,
    \item Projektdokumentation oder Deliverables,
    \item interne Kalkulation der GmbH,
    \item Vergleichsangebote, Marktpreisindikation oder interne Referenzsätze, soweit verfügbar.
\end{itemize}

\subsection{Fremdvergleichsergebnis}

Auf Grundlage des dargestellten Sachverhalts ist die Vergütung fremdüblich, sofern:

\begin{itemize}
    \item tatsächlich Leistungen erbracht werden,
    \item die Leistungen fachlich und zeitlich nachvollziehbar dokumentiert sind,
    \item die Vergütung dem Marktwert vergleichbarer Leistungen entspricht,
    \item die GmbH eine angemessene eigene Marge gegenüber dem Endkunden erzielt oder die Vergütung jedenfalls wirtschaftlich begründen kann,
    \item keine gesellschaftsrechtlich veranlassten Sondervorteile gewährt werden.
\end{itemize}

Unter diesen Voraussetzungen entspricht die Leistungsbeziehung dem Fremdvergleichsgrundsatz.

\section{Prüfung einer verdeckten Gewinnausschüttung}

\subsection{Prüfungsmaßstab}

Eine verdeckte Gewinnausschüttung kann vorliegen, wenn eine Kapitalgesellschaft ihrem Gesellschafter oder einer ihm nahestehenden Person einen Vorteil zuwendet, der durch das Gesellschaftsverhältnis veranlasst ist, sich auf den Gewinn der Gesellschaft auswirkt und einem fremden Dritten unter vergleichbaren Umständen nicht gewährt worden wäre.

Da Stefan Koch Gesellschafter der TEGRITY Safety Engineering GmbH ist, ist zu prüfen, ob die Zahlungen an ihn betrieblich veranlasst oder gesellschaftsrechtlich veranlasst sind.

\subsection{Betriebliche Veranlassung}

Für eine betriebliche Veranlassung sprechen:

\begin{itemize}
    \item die Leistungen werden tatsächlich erbracht,
    \item die Leistungen dienen Kundenprojekten der GmbH,
    \item die GmbH erhält einen wirtschaftlichen Nutzen,
    \item die Leistungen sind dokumentiert,
    \item die Vergütung orientiert sich am Fremdvergleich.
\end{itemize}

Die Zahlungen stehen damit dem Grunde nach im Zusammenhang mit dem operativen Geschäftsbetrieb der GmbH.

\subsection{Gesellschaftsrechtliche Veranlassung}

Gegen eine gesellschaftsrechtliche Veranlassung spricht:

\begin{itemize}
    \item Stefan Koch erhält keine Zahlung allein aufgrund seiner Gesellschafterstellung.
    \item Die Zahlung erfolgt für konkret nachweisbare Leistungen.
    \item Die GmbH hätte vergleichbare Leistungen auch von einem unabhängigen externen Berater beziehen können.
    \item Die Vergütung ist nicht unangemessen hoch.
    \item Es werden keine Gewinne ohne entsprechende Gegenleistung abgeschöpft.
\end{itemize}

\subsection{Besondere Risikopunkte}

Die folgenden Punkte wären vGA-relevant und sollten vermieden oder gesondert dokumentiert werden:

\begin{itemize}
    \item Vergütung ohne nachweisbare Leistung,
    \item überhöhte Vergütung im Vergleich zum Marktpreis,
    \item fehlende oder unklare Leistungsbeschreibung,
    \item Zahlungen für Tätigkeiten, die tatsächlich Geschäftsführungsfunktionen darstellen,
    \item nachträgliche oder rückwirkende Vergütungsvereinbarungen,
    \item fehlende Trennung zwischen Gesellschafterstellung und Dienstleistungsverhältnis,
    \item vollständige oder nahezu vollständige Abschöpfung der GmbH-Marge ohne wirtschaftliche Begründung.
\end{itemize}

\subsection{Prüfung dem Grunde nach}

Eine vGA dem Grunde nach liegt nicht vor, wenn die GmbH eine tatsächliche, werthaltige Gegenleistung erhält und die Beauftragung wirtschaftlich nachvollziehbar ist.

Dies ist der Fall, wenn die Leistungen für Kundenprojekte benötigt werden und die GmbH dadurch ihre eigenen vertraglichen Pflichten gegenüber Endkunden erfüllen kann.

\subsection{Prüfung der Höhe nach}

Eine vGA der Höhe nach liegt nicht vor, wenn die Vergütung marktüblich ist.

Die Vergütung ist insbesondere dann vertretbar, wenn:

\begin{itemize}
    \item sie sich an marktüblichen Tagessätzen, Stundensätzen oder Workpackage-Preisen orientiert,
    \item Qualifikation und Erfahrung des Leistungserbringers berücksichtigt werden,
    \item die GmbH eine wirtschaftlich nachvollziehbare Projektkalkulation vorlegen kann,
    \item die Vergütung nicht deutlich über dem liegt, was einem unabhängigen Dritten gezahlt worden wäre.
\end{itemize}

\subsection{vGA-Ergebnis}

Auf Grundlage des dargestellten Sachverhalts liegt keine verdeckte Gewinnausschüttung vor, sofern:

\begin{itemize}
    \item die Leistungen tatsächlich erbracht werden,
    \item die Leistungen nachvollziehbar dokumentiert sind,
    \item die Vergütung fremdüblich ist,
    \item keine Geschäftsführungs- oder Gesellschafterleistungen vergütet werden,
    \item die GmbH eine wirtschaftlich nachvollziehbare Gegenleistung erhält.
\end{itemize}

Die Zahlungen sind unter diesen Voraussetzungen betrieblich veranlasste Betriebsausgaben der GmbH und nicht durch das Gesellschaftsverhältnis veranlasst.

\section{Gesamtergebnis}

Die Prüfung führt zu folgendem Ergebnis:

\begin{enumerate}
    \item Eine DAC6-Meldepflicht besteht nicht, da kein relevantes Hallmark vorliegt und der Main-Benefit-Test nicht erfüllt ist.
    \item Die Vergütung ist fremdüblich, sofern sie dem Marktwert vergleichbarer Leistungen entspricht und durch Leistungsnachweise dokumentiert ist.
    \item Eine verdeckte Gewinnausschüttung liegt nicht vor, sofern tatsächliche Leistungen erbracht werden, die Vergütung angemessen ist und keine gesellschaftsrechtlich veranlassten Sondervorteile gewährt werden.
\end{enumerate}

\section{Empfohlene Anlagen}

Zur weiteren Absicherung sollten folgende Anlagen zu dieser Dokumentation genommen werden:

\begin{enumerate}
    \item Beratervertrag zwischen der TEGRITY Safety Engineering GmbH und Stefan Koch,
    \item Muster-Workpackage,
    \item konkrete Leistungsnachweise,
    \item Beispielrechnungen,
    \item interne Projektkalkulation der GmbH,
    \item kurze Marktpreis- oder Tagessatzbegründung,
    \item Nachweis der steuerlichen Ansässigkeit von Stefan Koch in Bulgarien, soweit verfügbar.
\end{enumerate}

\section{Schlussvermerk}

Diese Dokumentation wurde nach bestem Wissen auf Grundlage des dargestellten Sachverhalts erstellt. Sie ersetzt keine individuelle steuerliche Beratung durch einen Steuerberater oder Rechtsanwalt. Bei wesentlichen Änderungen des Sachverhalts ist eine erneute Prüfung erforderlich.

\vspace{2em}

\noindent
Ort, Datum: \rule{0.35\textwidth}{0.4pt}

\vspace{2em}

\noindent
Unterschrift / Verantwortlicher: \rule{0.45\textwidth}{0.4pt}

\end{document}